problem:  
Let \(M\) be a compact 7‑manifold equipped with a smooth \(G_2\)‑structure \(\varphi\) inducing a Riemannian metric \(g_\varphi\) and a corresponding spin structure. The Dirac operator on \((M,g_\varphi)\) is denoted \(D_\varphi\), and its square has discrete nonnegative spectrum \(\{ \lambda_i(\varphi) \}_{i \ge 0}\) with \(\lambda_0(\varphi)=0\) whenever \(\varphi\) is torsion‑free (since a parallel spinor exists).  

Define the *normalised first positive Dirac eigenvalue functional*  
\[
\mathcal{F}(\varphi) \;=\; \mathrm{Vol}(M,g_\varphi)^{1/7}\,\lambda_1(\varphi),
\]
where \(\lambda_1(\varphi)\) is the smallest strictly positive eigenvalue of \(D_\varphi^2\).  

**Conjecture.** Among all cohomologous \(G_2\)‑structures \(\varphi' \in [\varphi]\), the functional \(\mathcal{F}(\varphi')\) is minimized precisely at torsion‑free \(G_2\)‑structures.  

In other words, one speculates that the spectral quantity \(\mathcal{F}(\varphi)\) detects the vanishing of the intrinsic torsion: varying the \(G_2\)-structure within its cohomology class should increase (to first order) the lowest nonzero Dirac eigenvalue whenever torsion is introduced.  

Further, is it possible that the negative gradient flow of \(\mathcal{F}\) on the space of \(G_2\)‑structures, modulo diffeomorphisms, defines a new analytic flow whose critical points correspond exactly to torsion‑free \(G_2\)-structures (analogous to the Laplacian flow but spectrally guided)?  

Why is it a "good" problem:  
This problem is concrete while remaining profoundly open: it introduces a well-defined spectral functional associated with a \(G_2\)-structure and conjectures that torsion‑free structures are its minimizers. It ties together several deep areas—spin geometry (via Dirac operators), spectral analysis, and the differential geometry of special holonomy—and seeks a variational or analytic characterization of torsion‑free \(G_2\)-metrics purely through Dirac eigenvalues. Such a characterization would give a new spectral–geometric lens on exceptional holonomy, potentially leading to new analytic methods (e.g., gradient flows governed by Dirac spectra) for evolving arbitrary \(G_2\)-structures toward torsion‑free ones. The statement is simple enough to articulate clearly but conceptually deep, bridging global analysis and special holonomy theory—hallmarks of a “good” mathematical problem.